\documentclass[sigplan,nonacm]{acmart}
\makeatletter
\renewcommand\@formatdoi[1]{\ignorespaces}
\makeatother
\settopmatter{printfolios=false,printccs=false,printacmref=false}

\setcopyright{none}
\acmConference[]{}{}{}

\usepackage{enumitem}
\usepackage{mathtools}
\usepackage{amsmath}

\usepackage{listings}

\lstdefinelanguage{pcpp}{%
  language     = C++,
  morekeywords = {type,then,of,match,with,null},
  keywordstyle=\ttfamily\bfseries
}

\usepackage{natbib}
\usepackage{url}

\newcommand{\defeq}{\stackrel{\text{def}}{=}}
\newcommand{\setbuild}[2]{\{#1~|~#2\}}

% Syntax highlighting
\newdimen\zzlistingsize
\newdimen\zzlistingsizedefault
\zzlistingsizedefault=9pt
\newdimen\kwlistingsize
\kwlistingsize=9pt
\zzlistingsize=\zzlistingsizedefault
\gdef\lco{black}
\newcommand{\lcRed}{\gdef\lco{nred}}
\newcommand{\lcBlack}{\gdef\lco{black}}
\newcommand{\lcNormal}{\lcBlack}
\newcommand{\lcBlue}{\gdef\lco{nblue}}
\newcommand{\lcGreen}{\gdef\lco{ngreen}}
\newcommand{\lcPurple}{\gdef\lco{npurple}}
\newcommand{\Lstbasicstyle}{\fontsize{\zzlistingsize}{1.1\zzlistingsize}\ttfamily\color{\lco}}
\newcommand{\keywordstyle}{\fontsize{1.09\kwlistingsize}{\kwlistingsize}\normalfont\bf\color{\lco}}
\newlength{\zzlstwidth}
\settowidth{\zzlstwidth}{{\Lstbasicstyle~}}
\newcommand{\lcm}{\color{\lco}}

\lstset{
  basewidth=\zzlstwidth,
%  basicstyle=\Lstbasicstyle,
  basicstyle=\small\ttfamily,
%  identifierstyle=,
  numbers=left, numberstyle=\tiny, numbersep=5pt,
%  frame=single,
  keywordstyle=\keywordstyle,
  keywords={fun,cfun,sfun,read,case,of,type,functor,structure,signature,sig,struct,datatype,let,letpar,val,fn,in,end,as,select,mod,apply,sapply,capply,write},  %nil,cons,
  commentstyle=\rmfamily\slshape\color{\lco},
  morecomment=[n]{(*}{*)},
  lineskip={-0pt},
%  columns=[l]flexible,
  columns=[l]fullflexible,
  keepspaces=true,
  mathescape=true,
% NOTE: need TWO sets of braces around each definition below!
  literate={->}{{$\lcm\rightarrow$}}2
           {->^S}{{$\lcm\hspace{-1pt}\arrowsymbol{\S}\hspace{-1pt}$}}2
           {->^C}{{$\lcm\hspace{-1pt}\arrowsymbol{\C}\hspace{-1pt}$}}2
           {->^D}{{$\lcm\hspace{-1pt}\arrowsymbol{\delta}\hspace{-1pt}$}}2
           {->^Dh}{{$\lcm\hspace{-1pt}\arrowsymbol{\delta_h}\hspace{-1pt}$}}2
           {->^Dt}{{$\lcm\hspace{-1pt}\arrowsymbol{\delta_t}\hspace{-1pt}$}}2
%           {-c->}{${c\rightarrow}$}3
%           {-s->}{${s\rightarrow}$}3
%           {-c->}{$\rightarrow$}1
%           {-s->}{$\rightarrow$}1
%           {_}{{\lcm\_}}1
           {'a}{{$\alpha$}}1
           {'b}{{$\beta$}}1
           {'d}{{$\lcm\delta$}}1
}


\lstset{language=pcpp,numbers=none}

\newcommand{\fo}{\mathtt{f_1}}
\newcommand{\ft}{\mathtt{f_2}}
% A step number, flush right in a listing.
\newcommand{\stepc}[1]{\hfill{\rmfamily\slshape // (#1)}}
% A timeline: two columns of pseudocode, in the font of listings.
\newenvironment{timeline}[1]
  {\par\raggedright\noindent\small\ttfamily\begin{tabular}{@{}#1@{}}}
  {\end{tabular}\par}
\newcommand{\tlhead}[1]{{\rmfamily\itshape #1}}
\newcommand{\ind}{\hphantom{(0)~}}

\title{DAG calculus beside native fork join: one scheduler for counts and contexts}
% TODO: add co-author, and check affiliation.
\author{Mike Rainey}
\affiliation{
  \institution{Carnegie Mellon University}
  \city{Pittsburgh}
  \state{PA}
  \country{USA}
}
\email{me@mike-rainey.site}

\begin{document}

\begin{abstract}
Native fork join schedules fork-join programs by work stealing and
context switching, in a library, and makes a fork that is not stolen
cost little more than two calls.
%
But some parallel programs are not fork-join programs.
%
In a stencil computation, for example, each block of a grid waits for
its neighbors, and such dependencies are expressed directly in DAG
calculus, where a program builds a graph of \emph{vertices} and
\emph{edges} as it runs.
%
The two styles appear to need different schedulers: a fork keeps its
task in a stack frame and relies on a deque that holds only tasks of
its own stack, whereas a vertex lives on the heap and becomes ready
whenever its last predecessor finishes, on whichever worker that is.
%
In this short paper, we show that one scheduler serves both.
%
A vertex is a task with a count, and a stack that waits for a vertex
is itself a vertex, whose body is a saved context.
%
The scheduler of native fork join needs three changes: its loop takes
tasks from its own deque, a fork checks that the task it takes back is
its own, and stacks belong to workers rather than to tasks.
%
The fast path of a fork performs the same operations as before.
%
A companion implementation extends micro taskparts by about 190 lines
of C++, and includes two examples: a Gauss-Seidel stencil, and a
pipeline whose edges are added while it runs.
\end{abstract}

\maketitle

\section{Challenge: two styles, one scheduler}

Fork join suits computations that divide and conquer.
%
It suits others poorly.
%
Consider a Gauss-Seidel simulation of heat transfer on a grid, divided
into blocks.
%
At each time step, a block needs the new values of the blocks above it
and to its left, and the old values of the blocks below it and to its
right.
%
A fork-join program can run in parallel only blocks that it knows to be
independent, such as those on one diagonal, and must finish each such
group before it starts the next.
%
What the computation requires is weaker: a block may run a time step as
soon as its four neighbors allow it.
%
A pipeline is another example: its items are processed in parallel,
but each must wait for the one before it to be written out, which may
happen before, while, or after the item itself is processed.

DAG calculus~\cite{Dag_calculus_16} expresses such dependencies
directly.
%
A program creates \emph{vertices}, each of which holds a piece of code,
and adds \emph{edges} between them.
%
A vertex runs once it has been released and every vertex with an edge
to it has finished.
%
In an implementation, a vertex has an \emph{in-counter}, which counts
its unfinished predecessors, and an \emph{out-set}, which records its
successors.

What sets the calculus apart is that the number of edges that enter a
vertex and the number that leave it may both be decided as the program
runs.
%
Each of the usual styles fixes at least one of them.
%
In fork join, two branches enter a join, and each branch leaves for one
join.
%
In async-finish, any number of tasks may enter a finish, even tasks
spawned while the finish block runs, but each leaves for one finish.
%
With futures, any number of readers may leave a future, even readers
that arrive while the future is being computed, but each reader enters
with one future.
%
All three can be encoded in the calculus, and so can computations that
are none of them, such as the stencil.

We would like to offer both styles in one library, so that a vertex can
call \lstinline|fork2join|, and fork-join code can build a graph and
wait for it.
%
Encoding fork join in DAG calculus achieves that, but gives up what
native fork join~\cite{nativeforkjoin} provides: a fork that is not
stolen allocates nothing, and runs both branches as ordinary calls.
%
The opposite approach, two schedulers side by side, keeps the fast fork,
but leaves open how work moves between the two, and which of them an
idle worker should ask for work.
%
Our goal is one scheduler on which fork-join code runs as fast as it
does under native fork join alone.

At first sight, the scheduler of native fork join cannot host vertices,
because its safety argument relies on three facts that vertices break.
%
\begin{itemize}[leftmargin=*]
\item A worker's deque is empty whenever the worker is in its scheduler
  loop.
  But a vertex becomes ready when its last predecessor finishes, on a
  worker that may be in the middle of something else.
\item A deque holds only the tasks of forks in progress on one stack.
  So when a fork pops, it gets its own task or nothing.
  But a ready vertex belongs to no stack.
\item A task lives in a frame that outlives it, and owns the stack that
  it runs on.
  But a vertex lives on the heap, and may be freed, or run again, as
  soon as its successors learn that it has finished.
\end{itemize}

This paper shows that the scheduler needs only small changes to host
vertices.
%
Section~\ref{sec:nfj} recalls native fork join.
%
Section~\ref{sec:vertices} adds vertices, and
Section~\ref{sec:safety} revisits the safety argument.
%
Sections~\ref{sec:stencil} and~\ref{sec:pipeline} present two
examples, a stencil and a pipeline, and Section~\ref{sec:costs}
discusses costs and limits.
%
The appendices relate the design to DAG calculus and its original
implementation, and say what remains to be measured.

\section{Native fork join, in brief}
\label{sec:nfj}

We recall what the rest of the paper needs; the first paper of this
series~\cite{nativeforkjoin} has the details.

A \emph{context} holds the callee-saved registers, the stack pointer,
and a return address.
%
\lstinline|swap($k_1$, $k_2$)| saves the caller's context in $k_1$ and
resumes $k_2$, and \lstinline|restore($k$)| resumes $k$ without saving.
%
Any worker may resume a context, whichever worker saved it.

Each worker runs a scheduler loop on its thread's own stack, and client
code on separate \emph{task stacks}.
%
Each worker has a deque of tasks: the owner pushes and pops at one end,
and thieves steal the oldest task from the other.

A call to \lstinline|fork2join($\fo$, $\ft$)| keeps a \emph{join} and a
\emph{task} in its own frame.
%
The join holds a count of the branches that have yet to reach it,
initially 2, and a context.
%
The call pushes the task, runs $\fo$, and pops.
%
If the pop succeeds, the call runs $\ft$ and returns.
%
Otherwise the task was stolen: the call saves its context in the join
and switches to the scheduler loop.
%
A thief runs a stolen task on a fresh stack, and then switches to its
own loop.
%
The loop counts each branch that reaches a join, and the worker whose
decrement brings the count to zero resumes the context saved in the
join.

The pattern underneath is \emph{a count and a context}: a stack is
suspended until a count reaches zero, and a branch is counted only by a
scheduler loop, after the branch has left its stack.
%
Vertices fit the same pattern.

\section{Vertices}
\label{sec:vertices}

\begin{figure}
\hrule
\begin{lstlisting}
type join = {pending : atomic int, k : context}
type task =
  | fork of {body : () -> void, j : join*}
  | vertex of {body : () -> void,
      pending : atomic int, out : set of task*}
  | suspension of {k : context,
      pending : atomic int}
type worker = {k : context, dq : deque of task*,
  starting : task*, reached : join*,
  finished : task*, suspended : task*}

loop() -> void {  // on the worker's own stack
  while (!done) {
    j = me.reached; me.reached = null
    if (j != null &&
        fetch_and_add(j.pending, -1) == 1)
      { swap(me.k, j.k); continue }
    v = me.finished; me.finished = null
    if (v != null)
      for (u in close(v.out)) count_down(u)
    s = me.suspended; me.suspended = null
    if (s != null) count_down(s)
    t = pop()
    if (t == null) t = steal()
    if (t == null) continue
    if (t is suspension) swap(me.k, t.k)
    else { me.starting = t
      swap(me.k,
           new_context(new_stack(), enter)) } } }

count_down(v : task*) -> void {
  if (fetch_and_add(v.pending, -1) == 1)
    push(v) }  // v is ready

enter() -> void {  // first frame on a fresh stack
  t = me.starting; t.body()
  if (t is fork) {
    if (t.j == null) done = true  // root
    me.reached = t.j }
  else me.finished = t
  restore(me.k) }  // leave this stack
\end{lstlisting}
\hrule
\caption{The scheduler, with vertices}
\label{fig:scheduler}
\end{figure}

\begin{figure}
\hrule
\begin{lstlisting}[escapechar=@]
fork2join($\fo$, $\ft$ : () -> void) -> void {
  j = join{pending = 2}
  t = fork{body = $\ft$, j = &j}
  push(&t)@\stepc{1}@
  $\fo$()
  u = pop()@\stepc{2}@
  if (u == &t) { $\ft$(); return }@\stepc{3}@
  if (u != null) push(u)@\stepc{4}@
  me.reached = &j
  swap(j.k, me.k) }

new_vertex(f : () -> void) -> task* {
  return new vertex{body = f, pending = 1} }

new_edge(u, v : task*) -> bool {
  fetch_and_add(v.pending, 1)
  if (add(u.out, v)) return true
  fetch_and_add(v.pending, -1)  // u has finished
  return false }

release(v : task*) -> void { count_down(v) }

wait(v : task*) -> void {
  s = suspension{pending = 1}
  if (!new_edge(v, &s)) return
  me.suspended = &s
  swap(s.k, me.k) }
\end{lstlisting}
\hrule
\caption{Fork join, and the operations of DAG calculus}
\label{fig:ops}
\end{figure}

\paragraph{Three kinds of task}
%
A deque now holds three kinds of task
(Figure~\ref{fig:scheduler}).
%
A \emph{fork} is the second branch of a call to
\lstinline|fork2join|, as before.
%
A \emph{vertex} has a body, an in-counter, \lstinline|pending|, and an
out-set, \lstinline|out|.
%
A \emph{suspension} is the rest of a stack: it has a saved context
where the others have a body, and an in-counter too.
%
Whatever its kind, a task in a deque is ready to run and is not
running, and the worker that takes it out of the deque starts it,
once.
%
A fork or a vertex starts on a fresh stack, and a suspension starts
when its context is resumed.

\paragraph{The operations}
%
Figure~\ref{fig:ops} gives the operations of DAG calculus.
%
\lstinline|new_vertex| makes a vertex whose count is 1.
%
That one stands for the release that the vertex awaits, so the vertex
cannot become ready while its creator is still adding edges to it.
%
\lstinline|new_edge($u$, $v$)| makes $v$ wait for $u$: it increments
the count of $v$, and adds $v$ to the out-set of $u$.
%
If $u$ has already finished, its out-set is closed, the addition
fails, and \lstinline|new_edge| takes the increment back.
%
The caller must hold the count of $v$ above zero while it adds the
edge.
%
It does if $v$ has not been released, or if the caller is a vertex
that is running and has an edge to $v$.
%
The second case lets a vertex gain in-edges while its predecessors
run, which is what async-finish needs
(Appendix~\ref{sec:calculus}).
%
\lstinline|release($v$)| counts $v$ down.
%
In general, \lstinline|count_down| decrements the count of a task, and
pushes the task on the caller's deque if the count reaches zero.

\lstinline|wait($v$)| returns once $v$ has finished.
%
It makes a suspension in its own frame, adds an edge from $v$ to it,
saves its context in it, and switches to the scheduler loop.
%
The loop then counts the suspension down, which stands for its
release.
%
So a suspension is a vertex in all but its body, and to wait is to
turn the rest of one's stack into a vertex.
%
Any code that runs on a task stack may wait: the root, a branch of a
fork, or the body of a vertex.

\paragraph{The scheduler}
%
The loop treats a join as before.
%
It has three new duties.
%
When a vertex finishes, the loop closes its out-set and counts down
each task in it.
%
When a stack suspends in \lstinline|wait|, the loop counts down its
suspension.
%
And before it steals, the loop pops its own deque, which now may hold
tasks that the loop itself has just made ready.
%
As with joins, all of this happens in the loop, after the code
concerned has left its stack.

\paragraph{Fork join}
%
\lstinline|fork2join| differs from the original at steps 2 to 4.
%
The pop at step 2 may now return a task other than the one pushed at
step 1: a vertex that $\fo$ released, for instance.
%
So step 3 compares the task popped with its own.
%
If they differ, step 4 puts the other task back, and then suspends at
the join as before.
%
Section~\ref{sec:safety} explains why that is correct even if the
fork's own task is still in a deque.

\section{Why it is still safe}
\label{sec:safety}

The safety of native fork join rests on four
facts~\cite{nativeforkjoin}.
%
We restate each as it holds now.

\paragraph{Nothing is counted until its code has left its stack}
%
A join is counted down only by a scheduler loop, as before.
%
The same holds for the successors of a vertex: a vertex notifies them
from the loop, once its body has returned and the worker has left its
stack.
%
So when a successor runs, and frees the vertex or makes it ready
again, nothing of the vertex's previous run remains.
%
The loop reads the out-set with a single atomic exchange, which is its
last use of the vertex.
%
A suspension is released by the loop as well, after
\lstinline|swap| has saved its context, so no worker can resume a
context that is not yet saved.

\paragraph{No two workers run on the same stack}
%
A suspended stack is reachable only through a join or a suspension.
%
A join is resumed by the one worker that brings its count to zero.
%
A suspension is pushed on a deque by the one worker that brings its
count to zero, and resumed by the one worker that takes it out.
%
Its count cannot reach zero early, because only the waiting stack adds
edges to it, and the stack holds its release until it is saved.

\paragraph{A fork runs its own second branch, or leaves it to others}
%
The original argument was that the pop at step 2 returns the fork's own
task or nothing, because the deque held only tasks pushed on the
current stack.
%
That no longer holds, for two reasons.
%
A deque may hold ready vertices and suspensions.
%
And a stack may now suspend in \lstinline|wait| while tasks of its
forks are still in the deque, and resume on another worker.
%
The new argument is simpler, and does not depend on what the deque
holds.
%
If the pop returns the fork's own task, then no other worker has taken
it, and none will, so the fork may run $\ft$ itself.
%
Otherwise, the task was taken by another worker, or is still in some
deque.
%
In both cases it will run on a fresh stack and then reach the join, so
it is correct for the fork to suspend at the join, having put back
whatever it popped.
%
Because every scheduler loop now pops its own deque before it steals,
a task left in a deque is never stranded.

\paragraph{Only tasks that are ready and not running are in a deque}
%
A fork is pushed before it starts.
%
A vertex or a suspension is pushed when its count reaches zero.
%
Exactly one decrement brings a count to zero, and the count of a
vertex is not set again until the vertex has run and left its stack.
%
A join is never in a deque.

\paragraph{Whose stack?}
%
In the original, a stolen task records the stack it runs on, and the
frame that holds the task frees that stack.
%
A vertex has no such frame, and may be freed before its stack is.
%
So stacks now belong to workers.
%
A loop that starts a task takes a stack from its worker's cache, or
allocates one.
%
The first frame on the stack remembers the stack, and hands it to the
loop of whichever worker the task finishes on, which puts it in its
own cache.
%
The scheduler stores nothing in a task while the task runs, which is
what allows a vertex to be made ready again before the loop has
finished with its previous run.

\section{Example: a stencil}
\label{sec:stencil}

\begin{figure}
\hrule
\begin{lstlisting}
type block = vertex with {i, j, steps_left : int}

body(b : block*) -> void { update block (b.i, b.j) }

notify(b : block*) -> void {  // in the loop
  b.steps_left = b.steps_left - 1
  b.pending = number of neighbors of b
  count_down(below(b)); count_down(right(b))
  if (b.steps_left > 0)
    { count_down(above(b)); count_down(left(b)) }
  else if (b is the last block)
    for (u in close(b.out)) count_down(u) }

seidel(steps : int) -> void {
  for each block b {
    b.steps_left = steps
    b.pending = number of neighbors above
                and to the left of b }
  first.pending = 1; release(first)
  wait(last) }
\end{lstlisting}
\hrule
\caption{The Gauss-Seidel stencil.  Neighbors that do not exist are
  skipped.}
\label{fig:stencil}
\end{figure}

Our example is the Gauss-Seidel benchmark of the DAG calculus
paper~\cite{Dag_calculus_16}.
%
The grid is divided into square blocks.
%
One time step of a block replaces each of its cells, in place and in
row order, by the average of the cell and its four neighbors.
%
So time step $t$ of a block depends on step $t$ of the blocks above and
to the left, and on step $t-1$ of the blocks below and to the right.

Figure~\ref{fig:stencil} shows the program.
%
There is one vertex per block, which runs once per time step, so the
vertices are allocated once, in an array.
%
A block's edges are implicit: they go to its neighbors.
%
So a block replaces the out-set by code of its own, which the scheduler
loop runs in place of closing the out-set.
%
After each step, the block sets its count to the number of its
neighbors, for the next step, and then counts its neighbors down.
%
The block does all this in the loop, off its stack, which is why a
neighbor can make it ready again at once.
%
The caller releases the first block, and waits for the last, which
closes its out-set when it has run its final step.

The original implementation did not make a vertex of each block,
because starting a vertex was too expensive.
%
It kept the ready blocks in a queue inside one vertex, and split the
queue when another processor asked for work.
%
Here, a ready block is a pointer pushed on a deque, and the deque
already shares work on demand.

For comparison, the companion code has a fork-join version of the same
computation.
%
Block $(i, j)$ at step $t$ depends only on blocks whose
$i + j + 2t$ is smaller by one, so the blocks with the same
$i + j + 2t$ are independent.
%
The fork-join version runs these groups one after another, each as a
parallel loop.
%
Both versions compute, bit for bit, the grid that the serial program
computes.

\section{Example: a pipeline}
\label{sec:pipeline}

\begin{figure}
\hrule
\begin{lstlisting}
pipeline(K : int) -> void {
  i = 0
  while (more input) {
    x = read_chunk()  // stage 0
    dup = is_duplicate(x)  // stage 1
    if (i >= K) wait(v[i-K])  // throttle
    v[i] = new_vertex($\lambda$ () $\Rightarrow$ {
      if (!dup) y = compress(x)  // stage 2
      if (i > 0) wait(v[i-1])
      write(x, y, dup) })  // stage 3
    release(v[i]); i = i + 1 }
  if (i > 0) wait(v[i-1]) }
\end{lstlisting}
\hrule
\caption{The pipeline.  Inside the vertex, \lstinline|i|,
  \lstinline|x|, and \lstinline|dup| have the values they had when the
  vertex was created.}
\label{fig:pipeline}
\end{figure}

In the stencil, every edge is known before the computation starts.
%
Our second example needs edges that are not: the out-set of a vertex
grows after the vertex has started its work.
%
That is what a future requires, since a future may be forced at any
time, by any number of readers, and futures are one way to build
pipelines~\cite{blelloch1997pipelining}.

The example is a pipeline in the style of Cilk-P~\cite{lee2013piper}.
%
There, a loop processes a stream, each iteration passes through a
series of stages, and a stage of one iteration may wait for the same
stage of the previous iteration.
%
A stage that always waits is \emph{serial}: it runs for one iteration
at a time, in order.
%
A stage that never waits is \emph{parallel}.

Ours is the dedup benchmark, as Lee et al. pipeline it, much
simplified.
%
It compresses a stream of chunks, and replaces a chunk that it has
seen before by a reference.
%
Stage 0 reads a chunk.
%
Stage 1 hashes the chunk and looks the hash up in a table.
%
Stage 2 compresses the chunk, unless it is a duplicate.
%
Stage 3 writes the compressed chunk, or its hash.
%
Stage 2 is parallel, and the others are serial.
%
In our version, chunks have a fixed size, the hash is not
cryptographic, and the compressor is run-length encoding.

Figure~\ref{fig:pipeline} shows the program.
%
The loop runs the first two stages, which are both serial.
%
It makes one vertex per chunk, which runs stage 2, waits for the
vertex of the previous chunk, and runs stage 3.
%
So each vertex is a future that the next vertex forces.
%
The edge that this \lstinline|wait| adds is the dynamic one.
%
When stage 2 of a chunk ends, the previous vertex may be compressing,
suspended in its own \lstinline|wait|, writing, or finished.
%
In the first three cases, \lstinline|new_edge| adds a suspension to
the out-set of a vertex that is running, and in the last it finds the
out-set closed, and \lstinline|wait| returns at once.
%
The out-set settles the race between adding the edge and finishing the
vertex with one compare-and-swap.

The loop also waits.
%
Before it starts chunk $i$, it waits for chunk $i-K$, so at most $K$
chunks are in progress, and a loop that reads faster than the chunks
are compressed cannot fill memory with them.
%
This is the \emph{throttling} of Cilk-P, which likewise holds back
iteration $i+K$ until iteration $i$ has completed, and which used
$K = 4P$ for dedup on $P$ processors.
%
So a vertex in the middle of the stream may have two suspensions in its
out-set, that of the next vertex and that of the loop, added in either
order.
%
The loop also frees vertices: once vertex $i-K$ has finished, nothing
waits any longer for the vertices before it.

Cilk-P builds pipelines into its scheduler, with a construct for the
loop and one for each kind of stage.
%
Here, the pipeline is a client of the operations of
Figure~\ref{fig:ops}, and uses nothing else.
%
Where the stencil used a vertex that runs many times and whose edges
are implicit, the pipeline uses vertices that run once and whose edges
are added on the fly.
%
Its output is, byte for byte, that of the serial program.

\section{Costs and limits}
\label{sec:costs}

\paragraph{The cost of a fork}
%
A fork that is not stolen performs the same operations as before:
a push, a pop, and two calls.
%
The only difference is that it compares the result of the pop with the
address of its task, instead of with null.
%
A task no longer holds a context or a stack, so the frame of
\lstinline|fork2join| is smaller.
%
% TODO: cite a certified measurement of fib against the first paper's
% build, here and in Appendix sec:perf.

\paragraph{The cost of a vertex}
%
A vertex that is created costs one allocation, and each explicit edge
another.
%
A vertex that becomes ready costs a push, and one that starts costs a
stack from the worker's cache and two context switches: to the vertex,
and back to the loop.
%
A vertex whose edges are implicit, as in the stencil, allocates
nothing per run.

\paragraph{Mixing}
%
The two styles nest freely.
%
The body of a vertex runs on a task stack, so it may call
\lstinline|fork2join| and \lstinline|wait|.
%
If it suspends, its stack is resumed like any other, and the vertex
finishes when its body returns, on whichever worker that is.
%
Likewise, a branch of a fork may create vertices and wait for them.
%
A fork whose first branch releases vertices and does not wait for them
finds them in the deque at step 2.
%
It then takes the slow path, even though its task was not stolen: this
is correct, and costs a stack.

\paragraph{Limits}
%
The scheduler never frees a vertex, and the client must keep one alive
until it has finished and notified its successors, for instance by
waiting for a vertex downstream of it.
%
An edge may be added to a vertex only by code that holds the count of
the vertex above zero.
%
The code that a vertex runs in the scheduler loop, in place of closing
its out-set, runs on the worker's own stack, and so must neither fork
nor wait.
%
A deque no longer has a small bound: it holds the vertices that its
worker has made ready, as well as one task per fork in progress.
%
The consequences of native fork join for client code remain: code that
follows a fork or a wait may run on a different thread, and task stacks
have a fixed size.

\paragraph{High degrees}
%
The number of edges that enter or leave a vertex has no bound, and our
representation is the simplest one: an atomic counter, and a linked
list with a compare-and-swap at its head.
%
Two costs follow when a degree is large.
%
Many workers that add edges to one vertex, or count it down, contend
for one counter and one list head.
%
And one worker counts down all the successors of a vertex, one after
another, so a long out-set delays them.
%
Both have known remedies.
%
An in-counter can be a tree that grows as edges are added, so that
each operation costs amortized constant time, contention
included~\cite{acar2017dynsnzi}.
%
An out-set can be a tree as well, and its notification can be shared
among workers: the loop counts down part of the out-set, and hands the
rest to fresh vertices.
%
The original implementation of the calculus has
both~\cite{Dag_calculus_16}.
%
They would fit behind the two operations that touch the
representation, \lstinline|new_edge| and the code that notifies.
%
We have not built them.
%
For the out-set, we know of no program in practice whose out-degrees
are large enough to need one.
%
The in-counter is another matter.
%
Acar et al.~\cite{acar2017dynsnzi} measure a loop of asyncs with empty
bodies that all enter one finish, which is the pattern of a parallel
loop.
%
With a single atomic counter, such as ours, that loop ran fastest on
one core, and on every larger number of cores, up to forty, it ran
slower than with each of the tree counters they compared.
%
We expect our counter to behave likewise when many small vertices
enter one vertex, as in the async-finish of
Appendix~\ref{sec:calculus}, and we have not measured it.

\paragraph{Implementation}
%
The companion code is under \path{doc/dag-calculus} in the TaskPaRTS
repository~\cite{taskparts}.
%
It starts from the micro taskparts of the first
paper~\cite{nativeforkjoin}, and grows its scheduler from about 320
lines to about 510.
%
It includes the stencil in three versions (serial, fork join, and DAG
calculus), the pipeline in two (serial and DAG calculus), and a test
that mixes the two styles in both directions and includes an
async-finish.

\bibliographystyle{abbrv}
\bibliography{main}

\clearpage
\appendix

\section{Relation to DAG calculus}
\label{sec:calculus}

\paragraph{The calculus}
%
DAG calculus~\cite{Dag_calculus_16} has five operations.
%
\lstinline|newTd| creates a thread, which is our vertex.
%
\lstinline|newEdge| and \lstinline|release| are our
\lstinline|new_edge| and \lstinline|release|.
%
\lstinline|self| returns the running thread, and \lstinline|yield|
suspends it until the threads with edges to it have finished.
%
A thread typically creates threads, adds edges from them to itself,
releases them, and yields.

We have no \lstinline|self| and no \lstinline|yield|.
%
\lstinline|wait| covers that pattern, which is how the calculus
encodes fork join and the forcing of a future: the suspension that
\lstinline|wait| makes is the calculus's view of a thread that has
yielded, a vertex with edges entering it and the rest of the thread as
its body.
%
The encoding of async-finish needs one more vertex.
%
There, the threads spawned inside a finish block add edges to the
thread that will wait for them, and they do so while that thread is
still running the block.
%
A suspension cannot play that part, since it exists only once its
stack is about to be saved.
%
Instead, a finish block creates a vertex with an empty body to stand
for its end.
%
Each async creates a vertex, adds an edge from it to the end vertex,
and releases it.
%
The finish block then releases the end vertex and waits for it.
%
The end vertex gains in-edges while its predecessors run, and their
number is known only when the last of them has finished.
%
That is safe by the rule of Section~\ref{sec:vertices}: an async is
called by the finish block, which has yet to release the end vertex, or
by a vertex that has an edge to it.
%
A thread that yields in the calculus keeps its identity, and its
out-edges, across the yield.
%
Ours does too, in effect: a vertex finishes, and notifies its
out-set, only when its body returns, however many times its stack was
suspended in between.

\lstinline|fork2join| is the calculus's encoding of fork join,
specialized.
%
The encoding creates two threads, adds an edge from each to the
caller, releases them, and yields.
%
Native fork join creates one task, with no in-counter, because the
task is ready at once.
%
It runs the other branch in place.
%
And its join is an in-counter that can only be 2, 1, or 0, with the
caller's context beside it.

\paragraph{The original implementation}
%
The implementation that accompanies the calculus also switches
contexts: each vertex records a continuation, and the scheduler and
the vertices switch to each other as coroutines.
%
Three of its mechanisms have counterparts here.
%
It gives every new or running vertex an artificial in-edge, so that the
vertex cannot be scheduled while it is being set up, or while it is
still running.
%
Our release count is the first half of that edge.
%
We do not need the second half, because no edge enters a vertex once
its count has reached zero, and a vertex that runs again, as a block
of the stencil does, sets its count only after it has left its stack.
%
It adds an edge by incrementing the in-counter of the target first and
the out-set of the source second, and rolls back if the source has
finished, as \lstinline|new_edge| does.
%
And it notifies the out-set of a vertex from the scheduler, after the
vertex has finished.

It also lets each vertex choose a representation of its in-counter and
its out-set, from a single counter and a single pointer up to
concurrent trees for vertices of high degree.
%
Our vertex has one representation of each, which a vertex can replace,
as the blocks of the stencil do; Section~\ref{sec:costs} discusses
high degrees.

\section{What remains to be measured}
\label{sec:perf}

% TODO: this appendix is a plan. Replace it with certified Flexibench
% rounds, cited through \fbval, as in the first paper.

This draft cites no timings.
%
Five measurements are planned, each as a certified round.
%
\begin{itemize}[leftmargin=*]
\item \emph{The fork.}
  The running time of \lstinline|fib|, which does nothing but fork,
  with this paper's micro taskparts and with the first paper's, on one
  worker and on several.
  It shows whether hosting vertices costs anything on the fast path.
\item \emph{The stencil.}
  The three versions of the stencil, for a small grid with many time
  steps and a large grid with few, as in the DAG calculus paper, over
  a range of block sizes and numbers of workers.
  It shows what the weaker dependencies of the DAG version buy over
  the fork-join version.
\item \emph{The pipeline.}
  The two versions of the pipeline, over a range of chunk sizes,
  throttling limits, and numbers of workers.
  It shows what an edge added on the fly, and a suspension, cost when
  there is one of each per item.
\item \emph{The vertex.}
  The cost of creating, releasing, and waiting for a vertex whose body
  is empty, and of running a block of the stencil whose body is empty,
  beside the cost of a fork.
  It shows the overhead per vertex, and how far a vertex is from the
  goal of costing what a fork costs.
\item \emph{The in-counter.}
  The cost per async of an async-finish whose asyncs have empty or
  small bodies and all enter one finish, over a range of body sizes
  and numbers of workers.
  It shows whether a single atomic counter suffers from contention on
  this machine, as Section~\ref{sec:costs} expects.
\end{itemize}

\end{document}
